\documentclass[11pt,a4paper]{article}
\usepackage[utf8]{inputenc}
\usepackage[T1]{fontenc}
\usepackage[english]{babel}
\usepackage{geometry}
\geometry{a4paper, left=2.5cm, right=2.5cm, top=2.5cm, bottom=2.5cm}
\usepackage{mathptmx}
\usepackage{helvet}
\usepackage{amsmath}
\usepackage{titlesec}
\usepackage{booktabs}
\usepackage{tabularx}
\usepackage{xcolor}
\usepackage{authblk}
\usepackage{hyperref}
\usepackage{enumitem}
\usepackage{graphicx}
\usepackage{float}
\usepackage{setspace}
\usepackage{natbib}
\usepackage{longtable}
\usepackage{quoting}
\usepackage{multirow}
\usepackage{array}

\titleformat{\section}{\Large\bfseries\sffamily\color{black}}{\thesection}{1em}{}
\titleformat{\subsection}{\large\bfseries\sffamily\color{darkgray}}{\thesubsection}{1em}{}
\titleformat{\subsubsection}{\normalsize\bfseries\sffamily\color{darkgray}}{\thesubsubsection}{1em}{}

\hypersetup{
    pdftitle={An Integrated Multiaxial Model for Computer-Assisted Psychiatric Diagnosis},
    pdfauthor={Lukas Geiger},
    colorlinks=true,
    linkcolor=black,
    urlcolor=blue,
    citecolor=black
}

\onehalfspacing

\begin{document}

\title{\textbf{An Integrated Multiaxial Model for Computer-Assisted Psychiatric Diagnosis: Synthesis of DSM-5-TR, ICD-11, and ICF in a 6-Axis Expert System}\\[0.3em]
\large A Methods Paper}

\author{Lukas Geiger}
\date{February 2026 --- Version 2}

\maketitle

\section*{Abstract}

This paper presents a proposed 6-axis model for computer-assisted multiaxial psychiatric diagnosis that unifies categorical diagnostics (DSM-5-TR, ICD-11) with the functional classification of the ICF. The model addresses the clinical deficits created by the abolition of the multiaxial system in DSM-5 (2013): Axis~I (Mental Health Profiles) captures diagnoses across their temporal course with coverage analysis---defined as the systematic identification of diagnostically unexplained symptoms with a percentage-based metric---and a prioritized investigation plan; Axis~II (Biography and Development) integrates dimensional personality assessment (PID-5), formative experiences, and core conflicts; Axis~III (Medical Synopsis) mirrors Axis~I symmetrically with 13 sub-axes; Axis~IV (Environment and Functioning) combines WHODAS~2.0, ICF Core Sets, the Cultural Formulation Interview, and a structured treatment network; Axis~V (Integrated Condition Model) bridges case formulation with classification through an operationalized 3P/4P schema; Axis~VI (Evidence Collection and Clinical Safety) implements CAVE clinical alerts, longitudinal symptom tracking, and a contact log. The design philosophy rests on three principles: skeleton (structure without mandatory fields), living document, and report-as-snapshot. Each diagnosis receives PRO/CONTRA evidence evaluation with confidence estimation. HiTOP spectra are computed from Cross-Cutting screening results. The 6-step gatekeeper logic operationalizes the differential diagnosis sequence by First (2024). The coverage analysis extends existing decision-support methods (e.g., OPCRIT) by providing a transparent metric for clinical interpretability. The paper addresses dimensional integration (HiTOP/RDoC), ethics, a four-phase validation strategy, and HL7 FHIR interoperability. A Python reference implementation is publicly available at \url{https://github.com/lukisch/multiaxial-diagnostic-system}.

\textbf{Keywords:} Multiaxial Diagnosis, DSM-5-TR, ICD-11, Coverage Analysis, Expert System, HiTOP, HL7 FHIR, Case Formulation

\section{Introduction}

The abolition of the multiaxial diagnostic system in DSM-5 left a structural gap in psychiatric diagnostics. The original five-axis system introduced with DSM-III (1980) was eliminated due to poor inter-rater reliability of the GAF score, conceptual conflation of symptoms and functioning in Axis~V, inconsistent clinical use of Axis~IV, and the artificial boundary between Axes~I and~II. Available evidence suggests the clinical community lost more than it gained \citep{Kress2014, Probst2014}: the structured prompt for comprehensive biopsychosocial assessment was replaced by a flat listing that cannot replicate this function. Institutionally, the APA's Future DSM Strategic Planning Committee is developing a multi-domain diagnostic model, and \citet{ErlichFirst2025} explicitly propose returning to structured axial assessment.

\textbf{Existing multiaxial and integrative approaches.} The present model builds upon, but substantially extends, several foundational approaches to multiaxial and integrative psychiatric assessment. \citet{Williams1985} documented the origins and early critiques of the DSM-III multiaxial system, identifying both its clinical utility as a structured prompt for comprehensive assessment and its operational weaknesses---particularly the unreliable GAF score and the inconsistent application of the psychosocial stressors axis. The present model addresses these specific critiques: it replaces the GAF with validated instruments (WHODAS~2.0, ICF Core Sets), eliminates the artificial Axis~I/II boundary, and introduces quantitative coverage analysis as an objective complement to clinical judgment.

\citet{Bolton2008} argued that formulation---the individualized narrative account of a patient's difficulties---should complement, not replace, categorical diagnosis. Bolton's key insight was that classification systems capture \textit{what} a patient has, while formulation captures \textit{why} and \textit{how}. The present model operationalizes this insight through Axis~V (Integrated Condition Model), which implements a structured 3P/4P formulation schema with explicit links back to the diagnostic axes---bridging the gap between Bolton's narrative formulation and the structured classification that clinical systems require.

\citet{Wakefield1992} proposed the influential ``harmful dysfunction'' analysis, arguing that mental disorder is a condition in which an internal mechanism fails to perform its natural function \textit{and} the dysfunction causes harm to the individual. This conceptual framework informs the present model's approach to differential diagnosis: the 6-step gatekeeper logic \citep{First2024} systematically distinguishes normal variation from dysfunction by requiring evidence of both impaired mechanism \textit{and} clinically significant distress or functional impairment. The coverage analysis further operationalizes this by quantifying which symptoms are explained by established diagnoses and which remain unexplained---a direct operationalization of Wakefield's requirement to identify genuine dysfunction rather than mere statistical deviance.

The present model thus synthesizes these traditions: it retains the structural comprehensiveness of Williams's multiaxial vision, operationalizes Bolton's formulation imperative, and embeds Wakefield's conceptual rigor into a computationally tractable system---while adding proposed innovations (symmetric Axis~I/III architecture, quantitative coverage analysis, CAVE alerts) that aim to go beyond all three predecessors.

\section{Methods}

This paper follows a design science methodology for the development of a clinical decision-support artifact. The research process comprised four phases:

\textbf{Phase~1: Problem identification and requirements analysis.} A structured review of the literature on the abolition of the DSM-IV multiaxial system identified five specific clinical deficits: loss of the structured prompt for biopsychosocial assessment, poor GAF inter-rater reliability, inconsistent Axis~IV usage, the artificial Axis~I/II boundary, and the low adoption rate of replacement V/Z codes \citep{APA2013, Kress2014, Probst2014, ErlichFirst2025}. Requirements were derived from these deficits, from First's (2024) gold-standard differential diagnosis framework, and from current dimensional approaches (HiTOP, RDoC, ICD-11 dimensional personality model).

\textbf{Phase~2: Iterative design.} The 6-axis architecture was developed through iterative cycles of design, expert feedback, and refinement. Design decisions---such as the symmetric Axis~I/III structure, the coverage analysis innovation, and the CAVE alert system---were evaluated against the identified requirements. Each iteration was documented and versioned (V1--V5).

\textbf{Phase~3: Prototypical implementation.} A Python reference implementation (approximately 1,850 lines) using hierarchical state machines (\texttt{transitions} library) and a Streamlit interface was developed to demonstrate feasibility. The source code is publicly available at \url{https://github.com/lukisch/multiaxial-diagnostic-system}.

\textbf{Phase~4: AI-assisted development and review.} The manuscript and system design were developed with substantial AI assistance. Claude Opus~4.6 served as a drafting and synthesis tool; Gemini and Copilot were used as independent reviewers during revision. All AI contributions are documented in the AI Disclosure section (Disclosure Level~5). The scientific responsibility for all content rests with the human author.

\textbf{Scope and limitations of the present paper.} This is a methods and design paper. It presents the architecture, rationale, and prototypical implementation of the 6-axis model. Empirical validation (inter-rater reliability, clinical utility, convergent validity) is planned in four phases (Section~6) but has not yet been conducted. The coverage analysis thresholds (85\%/60\%) are proposed as initial values based on clinical reasoning and require empirical calibration in future studies.

\begin{table}[ht]
\centering
\caption{Comparison of Computer-Assisted Diagnostic Systems}
\begin{tabular}{lcccc}
\toprule
Feature & OPCRIT & CATEGO & MHIRA & Present Model \\
\midrule
Multiaxial structure & No & No & No & Yes (6 axes) \\
Dual coding (DSM/ICD) & Partial & ICD only & Flexible & Full \\
Coverage analysis & No & No & No & Yes \\
Dimensional integration & No & No & Partial & Yes (HiTOP) \\
FHIR interoperability & No & No & Partial & Planned \\
Empirical validation & Yes & Yes & Partial & Planned \\
\bottomrule
\end{tabular}
\label{tab:comparison}
\end{table}

\section{Architecture of the 6-Axis Model}

The model integrates categorical diagnostics (DSM-5-TR/ICD-11) into an epistemological total system. A central design principle is the \textbf{symmetric parallel structure} between Axis~I (psychological) and Axis~III (somatic): both axes feature analogous sub-axes (suspected diagnoses, treatment history, coverage analysis, investigation plan), enabling psychologists and physicians to work with identical structural tools in their respective domains. Axis~I tracks mental health profiles across ten sub-axes (Ia--Ij), with quantitative coverage analysis (Ii) preceding a prioritized 3-tier investigation plan (Ij, categorized as urgent/important/monitoring). Each diagnosis receives structured PRO/CONTRA evidence evaluation with numerical confidence estimation (0--100\%). Axis~II modernizes personality assessment through dimensional PID-5 trait profiling, supplemented by formative life experiences and core conflicts that feed into Axis~V case formulation. Axis~III encompasses 13 sub-axes: ten parallel the structure of Axis~I (IIIa--IIIj), with IIIb explicitly coding fully explanatory chronic conditions and IIIc explicitly coding contributing medical factors, plus three axis-specific sub-axes for integrated causal analysis (IIIk), genetic/family burden (IIIl), and structured medication history with efficacy ratings (IIIm). Axis~IV combines WHODAS~2.0, the German degree of disability (GdB, \textit{Grad der Behinderung}), ICF Core Sets, the Cultural Formulation Interview, and a structured contact persons and treatment network register. The GAF may optionally be retained for backward compatibility but is not a core component. Axis~V operationalizes case formulation through structured 3P/4P coding with explicit links to coverage analysis gaps. Axis~VI implements a central evidence matrix with clinical assessment fields, CAVE clinical alerts for cross-axis risk management (drug interactions, lab artifacts, contraindications, temporal misattributions, diagnostic limitations), longitudinal symptom tracking with therapy response documentation, and a contact and observation log documenting the diagnostic process itself. An exemplary multi-professional assignment model suggests psychologists for Axis~I, physicians for Axis~III, social workers for Axis~IV, and the interdisciplinary team for Axis~V synthesis; however, these assignments are non-mandatory and can be flexibly adapted to the clinical setting.

\section{The 6-Step Gatekeeper Logic}

The differential diagnosis sequence closely follows First's \citeyearpar{First2024} gold-standard framework: (1)~rule out malingering, (2)~rule out substance etiology, (3)~rule out medical conditions, (4)~determine primary disorders, (5)~differentiate adjustment disorders, (6)~establish boundary with no mental disorder. Step~6 is marked as ``conceptual'' because the threshold between a normal stress response and a mental disorder inherently requires clinical judgment, supported but not replaced by objective functional data.

\section{Cross-Cutting Screening}

The DSM-5 Cross-Cutting Symptom Measures (23~items across 13~domains for adults; 25~items across 12~domains for children/adolescents) form the screening backbone. The screening matrix now covers 19 domains including eating disorders (SCOFF, EDE-QS), sleep disorders (ISI), and impulse control, ensuring complete coverage across all 11 Hierarchical State Machine (HSM) disorder modules.

\section{Divergences, Coverage Analysis, and Dimensional Integration}

\subsection{Formal Definition of Coverage Analysis}

Let $S = \{s_1, \ldots, s_n\}$ be the set of observed symptoms and $D = \{d_1, \ldots, d_m\}$ the set of established diagnoses. For each symptom $s_i$, define the coverage function $c(s_i) = \max_{d_j \in D} w(s_i, d_j)$, where $w(s_i, d_j) \in [0, 1]$ represents the weight of symptom $s_i$ in diagnosis $d_j$ as defined by the diagnostic criteria. The total coverage is then $C_{\text{total}} = \frac{1}{n} \sum_{i=1}^{n} c(s_i)$. Symptom weights are derived from the diagnostic criteria of DSM-5-TR and ICD-11, with each criterion contributing equally within a diagnosis. Overlapping diagnoses are handled by assigning each symptom to the diagnosis providing maximum coverage.

\subsection{Classification System Divergences}

Five critical DSM-5-TR/ICD-11 divergences require dual-system encoding, including Prolonged Grief Disorder (newly introduced in DSM-5-TR; present in ICD-11 as Prolonged Grief Disorder, 6B42). The coverage analysis is implemented symmetrically in Axis~I (Ii) and Axis~III (IIIi): each symptom receives a coverage percentage, classified as fully covered ($\geq$85\%), partially covered (60--84\%), or insufficiently covered ($<$60\%), yielding a total coverage score. Dimensional integration via automatic HiTOP spectra computation from Cross-Cutting domain scores and RDoC research annotation provides complementary perspectives to categorical diagnosis.

\section{Ethics, Validation, and Interoperability}

\subsection{Ethics and Data Protection}
The system processes sensitive psychiatric data and must comply with GDPR (EU 2016/679) and applicable national data protection regulations. All patient data are stored locally with AES-256 encryption. The system is designed as a clinical decision support tool, not as an autonomous diagnostic agent---final diagnostic authority remains with the clinician. Algorithmic bias monitoring is planned through regular audits of coverage analysis outputs across demographic subgroups.

\subsection{Planned Validation Strategy}
Validation is planned in four phases: (1) Expert panel evaluation of axis structure and completeness (N = 15--20 clinicians), (2) Pilot implementation in a clinical setting with usability assessment (System Usability Scale), (3) Criterion validity study comparing coverage analysis outputs with independent clinical diagnoses (N $\geq$ 200), and (4) Multi-site field trial assessing inter-rater reliability and diagnostic concordance.

\subsection{Interoperability}
The system targets HL7 FHIR R4 compatibility. Axis I diagnoses map to FHIR \texttt{Condition} resources, Axis III to \texttt{Observation} resources, and Axis IV functional assessments to \texttt{ClinicalImpression} resources. Full FHIR mapping specifications are planned for a subsequent technical report.

\section{Limitations}

This paper has several significant limitations that must be acknowledged.

\textbf{No empirical validation.} The 6-axis model has not been tested with real patients or clinicians. The coverage analysis thresholds (85\%/60\%), the CAVE alert triggers, and the multi-professional assignment model are all based on theoretical reasoning and clinical experience, not on empirical data. Until inter-rater reliability studies, clinical utility assessments, and convergent validity analyses are conducted (as outlined in the four-phase validation roadmap, Section~7), the model's practical value remains undemonstrated.

\textbf{Coverage analysis requires empirical calibration.} The coverage analysis---a proposed approach for which we are not aware of directly comparable implementations, though related approaches exist---is formally defined in Section~6 but has not been empirically calibrated. The coverage analysis is intuitively appealing but raises several unresolved questions: What are typical coverage scores in different diagnostic populations? How sensitive is the total coverage score to individual symptom ratings? The proposed thresholds (85\%/60\%) require empirical validation through the planned field trials before clinical adoption.

\textbf{Complexity may limit adoption.} The model's comprehensiveness (6 axes, 36+ sub-axes, 11 disorder modules) may constitute a barrier to clinical adoption. The historical failure of DSM-IV's multiaxial system was partly due to inconsistent clinical use---a simpler system has a better chance of consistent application. The ``skeleton principle'' (structure without mandatory fields) mitigates but does not eliminate this risk.

\textbf{Limited literature base.} As a design paper, this work draws primarily on diagnostic classification manuals and a small number of foundational papers. A systematic review of all existing computer-assisted diagnostic systems (e.g., OPCRIT, DTREE, CATEGO) and their empirical track records would strengthen the rationale for the specific design choices made here.

\textbf{AI-heavy development process.} The substantial AI involvement in both system design and manuscript preparation (Disclosure Level~5, as defined in the AI Disclosure section) raises questions about the reproducibility and independent verifiability of design decisions. While the source code is publicly available, the iterative design process---including expert feedback cycles---is not fully documented in a way that permits independent replication.

\section{Conclusion}

The proposed 6-axis system addresses a structural gap in psychiatric diagnostics that has persisted since the abolition of the DSM-IV multiaxial system. For clinical practice, the model's primary contribution lies in restoring---and substantially extending---the structured prompt for comprehensive biopsychosocial assessment. The symmetric Axis~I/III architecture enables genuine multi-professional collaboration, the coverage analysis provides clinicians with a transparent, quantitative metric for diagnostic completeness, and the CAVE alert system addresses the growing challenge of cross-axis risk management in complex cases. The PRO/CONTRA evidence evaluation with confidence estimation operationalizes diagnostic transparency in a way that supports both clinical accountability and training.

The model's architecture is consistent with the direction in which major classification systems are evolving. The APA's Future DSM Strategic Planning Committee is exploring dimensionality, biomarkers, and social determinants---all of which the present model already accommodates through HiTOP integration, RDoC annotation readiness, and the structured Axis~IV. The convergence with the explicit call by \citet{ErlichFirst2025} for a return to structured axial assessment supports the timeliness of this contribution. At the same time, the model extends beyond what institutional revision alone is likely to achieve, particularly through the coverage analysis and the operationalized case formulation (Axis~V) that bridges classification and treatment planning.

Concrete next steps include: (1) conducting the Phase~1 expert panel review to refine the axis structure, (2) executing the inter-rater reliability study with standardized case vignettes (Phase~2), (3) developing the HL7 FHIR mapping specification for interoperability with existing health IT systems, and (4) publishing the formal mathematical specification of the coverage analysis with worked clinical examples. The complete source code is publicly available to facilitate independent evaluation and replication.


% ===================================================================
% LITERATURVERZEICHNIS
% ===================================================================
\newpage
\begin{thebibliography}{99}

\bibitem[Aboraya et al.(2006)]{Aboraya2006}
Aboraya, A., Rankin, E., France, C. et al. (2006).
\newblock The reliability of psychiatric diagnosis revisited.
\newblock \textit{Psychiatry}, 3(1), 41--50.

\bibitem[Allison et al.(2012)]{Allison2012}
Allison, C., Auyeung, B. \& Baron-Cohen, S. (2012).
\newblock Toward brief ``red flags'' for autism screening: The Short Autism Spectrum Quotient and the Short Quantitative Checklist in 1,000 cases and 3,000 controls.
\newblock \textit{Journal of the American Academy of Child \& Adolescent Psychiatry}, 51(2), 202--212.

\bibitem[APA(2013)]{APA2013}
American Psychiatric Association (2013).
\newblock \textit{Diagnostic and Statistical Manual of Mental Disorders} (5th ed.).
\newblock Arlington, VA: American Psychiatric Publishing.

\bibitem[APA(2022)]{APA2022}
American Psychiatric Association (2022).
\newblock \textit{Diagnostic and Statistical Manual of Mental Disorders, Text Revision} (DSM-5-TR).
\newblock Washington, DC: American Psychiatric Publishing.

\bibitem[Ayuso-Mateos et al.(2010)]{AyusoMateos2010}
Ayuso-Mateos, J.\,L., Miret, M., Caballero, F.\,F. et al. (2010).
\newblock Multi-country evaluation of ICF Core Sets for depression.
\newblock \textit{BMC Psychiatry}, 10, 44.

\bibitem[Bach et al.(2022)]{Bach2022}
Bach, B., Kramer, U., Doering, S. et al. (2022).
\newblock The ICD-11 classification of personality disorders: A European perspective on challenges and opportunities.
\newblock \textit{Borderline Personality Disorder and Emotion Dysregulation}, 9(1), 12.

\bibitem[Barkham et al.(2023)]{Barkham2023}
Barkham, M., De Jong, K., Delgadillo, J. \& Lutz, W. (2023).
\newblock Routine outcome monitoring (ROM) and feedback.
\newblock \textit{Psychotherapy Research}, 33(4), 407--413.

\bibitem[Blevins et al.(2015)]{Blevins2015}
Blevins, C.\,A., Weathers, F.\,W., Davis, M.\,T. et al. (2015).
\newblock The Posttraumatic Stress Disorder Checklist for DSM-5 (PCL-5).
\newblock \textit{Journal of Traumatic Stress}, 28(6), 489--498.

\bibitem[Bolton(2008)]{Bolton2008}
Bolton, D. (2008).
\newblock \textit{What is Mental Disorder? An Essay in Philosophy, Science, and Values}.
\newblock Oxford: Oxford University Press.

\bibitem[Carlson \& Putnam(1993)]{CarlsonPutnam1993}
Carlson, E.\,B. \& Putnam, F.\,W. (1993).
\newblock An update on the Dissociative Experiences Scale.
\newblock \textit{Dissociation}, 6(1), 16--27.

\bibitem[Cloitre et al.(2018)]{Cloitre2018}
Cloitre, M., Shevlin, M., Brewin, C.\,R. et al. (2018).
\newblock The International Trauma Questionnaire: Development of a self-report measure of ICD-11 PTSD and Complex PTSD.
\newblock \textit{Acta Psychiatrica Scandinavica}, 138(6), 536--546.

\bibitem[Erlich et al.(2025)]{ErlichFirst2025}
Erlich, M.\,D., First, M.\,B., Talley, R.\,M. et al. (2025).
\newblock Integrating social determinants of health into clinical formulations: The need for a biaxial system in the DSM and psychiatry.
\newblock \textit{Psychiatric Services}, 76(10), 911--914.

\bibitem[First(2024)]{First2024}
First, M.\,B. (2024).
\newblock \textit{DSM-5-TR Handbook of Differential Diagnosis}.
\newblock Washington, DC: American Psychiatric Publishing.

\bibitem[Foa et al.(2002)]{Foa2002}
Foa, E.\,B., Huppert, J.\,D., Leiberg, S. et al. (2002).
\newblock The Obsessive-Compulsive Inventory: Development and validation of a short version.
\newblock \textit{Psychological Assessment}, 14(4), 485--496.

\bibitem[Gspandl et al.(2018)]{Gspandl2018}
Gspandl, S., Peirson, R.\,P., Nahhas, R.\,W., Skale, T.\,G. \& Lehrer, D.\,S. (2018).
\newblock Comparing Global Assessment of Functioning (GAF) and World Health Organization Disability Assessment Schedule (WHODAS) 2.0 in schizophrenia.
\newblock \textit{Psychiatry Research}, 259, 251--253.

\bibitem[Hevner et al.(2004)]{Hevner2004}
Hevner, A.\,R., March, S.\,T., Park, J. \& Ram, S. (2004).
\newblock Design science in information systems research.
\newblock \textit{MIS Quarterly}, 28(1), 75--105.

\bibitem[Huber et al.(2022)]{Huber2022}
Huber, C.\,G., Aebi, B.\,J., Engeler, S.\,L. et al. (2022).
\newblock Adoption of routine outcome monitoring in mental health settings.
\newblock \textit{BMC Health Services Research}, 22, 1035.

\bibitem[Insel et al.(2010)]{Insel2010}
Insel, T., Cuthbert, B., Garvey, M. et al. (2010).
\newblock Research Domain Criteria (RDoC): Toward a new classification framework for research on mental disorders.
\newblock \textit{American Journal of Psychiatry}, 167(7), 748--751.

\bibitem[Ising et al.(2012)]{Ising2012}
Ising, H.\,K., Veling, W., Loewy, R.\,L. et al. (2012).
\newblock The validity of the 16-item version of the Prodromal Questionnaire (PQ-16) to screen for ultra high risk of developing psychosis.
\newblock \textit{Schizophrenia Research}, 138(2--3), 207--212.

\bibitem[Keeley et al.(2016)]{Keeley2016}
Keeley, J.\,W., Reed, G.\,M., Roberts, M.\,C. et al. (2016).
\newblock Developing a science of clinical utility in diagnostic classification systems: Field study strategies for ICD-11 mental and behavioural disorders.
\newblock \textit{American Psychologist}, 71(1), 3--16.

\bibitem[Kessler et al.(2005)]{Kessler2005}
Kessler, R.\,C., Adler, L., Ames, M. et al. (2005).
\newblock The World Health Organization Adult ADHD Self-Report Scale (ASRS).
\newblock \textit{Psychological Medicine}, 35(2), 245--256.

\bibitem[Kotov et al.(2017)]{Kotov2017}
Kotov, R., Krueger, R.\,F., Watson, D. et al. (2017).
\newblock The Hierarchical Taxonomy of Psychopathology (HiTOP): A dimensional alternative to traditional nosologies.
\newblock \textit{Journal of Abnormal Psychology}, 126(4), 454--477.

\bibitem[Kotov et al.(2021)]{Kotov2021}
Kotov, R., Krueger, R.\,F., Watson, D. et al. (2021).
\newblock The Hierarchical Taxonomy of Psychopathology (HiTOP): A quantitative nosology based on consensus of evidence.
\newblock \textit{Annual Review of Clinical Psychology}, 17, 83--108.

\bibitem[Kotov et al.(2022)]{Kotov2022}
Kotov, R., Krueger, R.\,F. \& Watson, D. (2022).
\newblock The Hierarchical Taxonomy of Psychopathology (HiTOP) in psychiatric practice and research.
\newblock \textit{Psychological Medicine}, 52(9), 1666--1678.

\bibitem[Kress et al.(2014)]{Kress2014}
Kress, V.\,E., Barrio Minton, C.\,A., Adamson, N.\,A., Paylo, M.\,J. \& Pope, V. (2014).
\newblock The removal of the multiaxial system in the DSM-5: Implications and practice suggestions for counselors.
\newblock \textit{The Professional Counselor}, 4(3), 191--201.

\bibitem[Levis et al.(2019)]{Levis2019}
Levis, B., Benedetti, A. \& Thombs, B.\,D. (2019).
\newblock Accuracy of Patient Health Questionnaire-9 (PHQ-9) for screening to detect major depression.
\newblock \textit{BMJ}, 365, l1476.

\bibitem[Linden \& Baron(2005)]{LindenBaron2005}
Linden, M. \& Baron, S. (2005).
\newblock Das Mini-ICF-APP: Mini-ICF-Rating f\"ur Aktivit\"ats- und Partizipationsst\"orungen bei psychischen Erkrankungen.
\newblock \textit{Die Rehabilitation}, 44(3), 153--159.

\bibitem[Gierk et al.(2014)]{Gierk2014}
Gierk, B., Kohlmann, S., Kroenke, K., Spangenberg, L., Zenger, M., Br\"ahler, E. \& L\"owe, B. (2014).
\newblock The Somatic Symptom Scale-8 (SSS-8): A brief measure of somatic symptom burden.
\newblock \textit{JAMA Internal Medicine}, 174(3), 399--407.

\bibitem[Luciano et al.(2010)]{Luciano2010}
Luciano, J.\,V., Ayuso-Mateos, J.\,L., Aguado, J. et al. (2010).
\newblock The 12-item World Health Organization Disability Assessment Schedule II (WHO-DAS II): A nonparametric item response analysis.
\newblock \textit{BMC Medical Research Methodology}, 10, 45.

\bibitem[Morgan et al.(1999)]{Morgan1999}
Morgan, J.\,F., Reid, F. \& Lacey, J.\,H. (1999).
\newblock The SCOFF questionnaire: Assessment of a new screening tool for eating disorders.
\newblock \textit{BMJ}, 319(7223), 1467--1468.

\bibitem[Morin et al.(2011)]{Morin2011}
Morin, C.\,M., Belleville, G., B\'elanger, L. \& Ivers, H. (2011).
\newblock The Insomnia Severity Index: Psychometric indicators to detect insomnia cases and evaluate treatment response.
\newblock \textit{Sleep}, 34(5), 601--608.

\bibitem[Narrow et al.(2013)]{Narrow2013}
Narrow, W.\,E., Clarke, D.\,E., Kuramoto, S.\,J. et al. (2013).
\newblock DSM-5 field trials in the United States and Canada, Part III.
\newblock \textit{American Journal of Psychiatry}, 170(1), 71--82.

\bibitem[Nordgaard et al.(2020)]{Nordgaard2020}
Nordgaard, J., Henriksen, M.\,G., Berge, J. \& Nilsson, L.\,S. (2020).
\newblock First-admission patients in a psychiatric acute-care setting: Nonspecificity of psychiatric symptoms and diagnostic instability.
\newblock \textit{Nordic Journal of Psychiatry}, 74(8), 581--586.

\bibitem[Oltmanns \& Widiger(2018)]{OltmannsWidiger2018}
Oltmanns, J.\,R. \& Widiger, T.\,A. (2018).
\newblock A self-report measure for the ICD-11 dimensional trait model proposal: The Personality Inventory for ICD-11.
\newblock \textit{Psychological Assessment}, 30(2), 154--169.

\bibitem[Owen(2023)]{Owen2023}
Owen, G. (2023).
\newblock What is formulation in psychiatry?
\newblock \textit{Psychological Medicine}, 53(5), 1700--1707. DOI: 10.1017/S0033291723000016.

\bibitem[Posner et al.(2011)]{Posner2011}
Posner, K., Brown, G.\,K., Stanley, B. et al. (2011).
\newblock The Columbia-Suicide Severity Rating Scale (C-SSRS).
\newblock \textit{American Journal of Psychiatry}, 168(12), 1266--1277.

\bibitem[Probst(2014)]{Probst2014}
Probst, B. (2014).
\newblock The life and death of Axis IV: Caught in the quest for a theory-free diagnostic system.
\newblock \textit{Research on Social Work Practice}, 24(1), 123--131.

\bibitem[Rosenman et al.(2003)]{Rosenman2003}
Rosenman, S.\,J., Korten, A.\,E., Medway, J. \& Evans, M. (2003).
\newblock Dimensional vs.\ categorical diagnosis in psychosis.
\newblock \textit{Acta Psychiatrica Scandinavica}, 107(5), 378--384.

\bibitem[Saunders et al.(1993)]{Saunders1993}
Saunders, J.\,B., Aasland, O.\,G., Babor, T.\,F. et al. (1993).
\newblock Development of the Alcohol Use Disorders Identification Test (AUDIT).
\newblock \textit{Addiction}, 88(6), 791--804.

\bibitem[Spitzer et al.(2006)]{Spitzer2006}
Spitzer, R.\,L., Kroenke, K., Williams, J.\,B.\,W. \& L\"owe, B. (2006).
\newblock A brief measure for assessing Generalized Anxiety Disorder: The GAD-7.
\newblock \textit{Archives of Internal Medicine}, 166(10), 1092--1097.

\bibitem[Stasiak et al.(2023)]{Stasiak2023}
Stasiak, K., Merry, S., Lignos, A. et al. (2023).
\newblock Mental Health Information Reporting Assistant (MHIRA)---an open-source software facilitating evidence-based assessment for clinical services.
\newblock \textit{BMC Psychiatry}, 23, 706.

\bibitem[Wakefield(1992)]{Wakefield1992}
Wakefield, J.\,C. (1992).
\newblock The concept of mental disorder: On the boundary between biological facts and social values.
\newblock \textit{American Psychologist}, 47(3), 373--388.

\bibitem[WHO(2010)]{WHO2010}
World Health Organization (2010).
\newblock \textit{WHODAS 2.0: WHO Disability Assessment Schedule 2.0}.
\newblock Geneva: WHO.

\bibitem[WHO(2019)]{WHO2019}
World Health Organization (2019).
\newblock \textit{International Classification of Diseases, 11th Revision} (ICD-11).
\newblock Geneva: WHO.

\bibitem[Williams(1985)]{Williams1985}
Williams, J.\,B.\,W. (1985).
\newblock The multiaxial system of DSM-III: Where did it come from and where should it go? I.\ Its origins and critiques.
\newblock \textit{Archives of General Psychiatry}, 42(2), 175--180.

\bibitem[Zimmermann et al.(2021)]{Zimmermann2021}
Zimmermann, J., Kerber, A., Rek, K. et al. (2021).
\newblock Psychometric properties of the independent 36-item PID5BF+M for ICD-11 in the Czech-speaking community sample.
\newblock \textit{Frontiers in Psychiatry}, 12, 643270.

\bibitem[McGuffin et al.(1991)]{McGuffin1991}
McGuffin, P., Farmer, A. \& Harvey, I. (1991).
\newblock A polydiagnostic application of operational criteria in studies of psychotic illness: Development and reliability of the OPCRIT system.
\newblock \textit{Archives of General Psychiatry}, 48(8), 764--770.

\bibitem[Wing et al.(1974)]{Wing1974}
Wing, J.\,K., Cooper, J.\,E. \& Sartorius, N. (1974).
\newblock \textit{Measurement and Classification of Psychiatric Symptoms: An Instruction Manual for the PSE and CATEGO Program}.
\newblock Cambridge: Cambridge University Press.

\bibitem[Peffers et al.(2007)]{Peffers2007}
Peffers, K., Tuunanen, T., Rothenberger, M.\,A. \& Chatterjee, S. (2007).
\newblock A design science research methodology for information systems research.
\newblock \textit{Journal of Management Information Systems}, 24(3), 45--77.

\bibitem[Krueger et al.(2018)]{Krueger2018}
Krueger, R.\,F., Kotov, R., Watson, D. et al. (2018).
\newblock Progress in achieving quantitative classification of psychopathology.
\newblock \textit{World Psychiatry}, 17(3), 282--293.

\bibitem[Regier et al.(2013)]{Regier2013}
Regier, D.\,A., Narrow, W.\,E., Clarke, D.\,E. et al. (2013).
\newblock DSM-5 field trials in the United States and Canada, Part II: Test-retest reliability of selected categorical diagnoses.
\newblock \textit{American Journal of Psychiatry}, 170(1), 59--70.

\end{thebibliography}

\end{document}
